\documentclass[a4paper,12pt]{article}

% Document Setup
\usepackage[utf8]{inputenc}
\usepackage[T1]{fontenc}
\usepackage[french]{babel}
\usepackage[margin=2.5cm]{geometry}
\usepackage{enumitem}
\usepackage{tikz}
\usetikzlibrary{shapes.cloud, shadows, positioning}
\usepackage{tcolorbox}
\usepackage{fontawesome5}
\usepackage{xcolor}
\usepackage{sectsty}
\usepackage{hyperref}
\usepackage{listings}
\usepackage{xcolor}

% Color Palette - Deep Slate Blue and Soft Teal/Cyan
\definecolor{primaryBlue}{RGB}{59, 76, 103}    % Deep Slate Blue
\definecolor{accentTeal}{RGB}{77, 166, 166}     % Soft Teal/Cyan
\definecolor{codeBg}{RGB}{248, 249, 250}
\definecolor{codeBorder}{RGB}{59, 76, 103}

% Section styling
\sectionfont{\color{primaryBlue}\normalfont}
\subsectionfont{\color{accentTeal}\normalfont}

% Code listing style
\lstset{
  backgroundcolor=\color{codeBg},
  basicstyle=\ttfamily\small,
  breaklines=true,
  frame=single,
  frameround=tttt,
  rulecolor=\color{codeBorder},
  keywordstyle=\color{primaryBlue},
  commentstyle=\color{accentTeal},
  stringstyle=\color{primaryBlue!70!black},
}

\pagestyle{empty}

\title{\color{primaryBlue}\Huge\bfseries Agent Extracteur : Architecture et Documentation Technique}
\date{\today}

\begin{document}

\maketitle

% TikZ Diagram - Project Architecture
\begin{center}
\begin{tikzpicture}[
  node distance = 2.2cm,
  component/.style = {rectangle, draw=primaryBlue, fill=primaryBlue!10, 
                      rounded corners, minimum width=3.2cm, minimum height=1.3cm,
                      align=center, font=\small, drop shadow},
  module/.style = {rectangle, draw=accentTeal, fill=accentTeal!15, 
                   thick, rounded corners, minimum width=2.8cm, minimum height=1.1cm,
                   align=center, font=\footnotesize},
  interface/.style = {cloud, draw=primaryBlue, fill=primaryBlue!20, 
                      cloud puffs=10, aspect=2, align=center}
]

% External interfaces
\node[interface] (sources) {Sources externes\\Syslog, Réseau, SIEM};

% Main agent
\node[component, below=of sources] (agent) {\faIcon{robot} SecurityAgent};

% Internal modules
\node[module, below left=1.2cm and 2cm of agent] (syslog) {Syslog Receiver};
\node[module, below=of agent] (network) {Network Capture};
\node[module, below right=1.2cm and 2cm of agent] (siem) {SIEM Collector};

% Output
\node[interface, below=1.8cm of agent] (memory) {Shared Memory\\(data/processed)};

% Arrows
\draw[->, thick, primaryBlue] (sources) -- (agent);
\draw[->, thick, accentTeal] (agent) -- (syslog);
\draw[->, thick, accentTeal] (agent) -- (network);
\draw[->, thick, accentTeal] (agent) -- (siem);
\draw[->, thick, primaryBlue] (syslog) -- (memory);
\draw[->, thick, primaryBlue] (network) -- (memory);
\draw[->, thick, primaryBlue] (siem) -- (memory);

\end{tikzpicture}
\end{center}

\section*{Description générale du projet}

L'\textbf{Agent Extracteur} est une composante centrale du système multi-agents de cybersécurité, conçue pour l'\textbf{ingestion unifiée} et la \textbf{pré-analyse} des données de sécurité provenant de multiples sources heterogènes.

\begin{tcolorbox}[
  enhanced,
  colback=codeBg,
  colframe=primaryBlue,
  left=8pt,
  right=8pt,
  top=8pt,
  bottom=8pt,
  boxrule=2pt,
  arc=4pt,
  before skip=12pt,
  after skip=12pt
]
\textbf{Objectif principal} : Assurer une collecte fiable, asynchrone et sans perte de données de sécurité depuis les sources de logs, le trafic réseau et les alertes SIEM, avec un mécanisme d'arrêt gracieux et une journalisation complète des erreurs.
\end{tcolorbox}

\section*{Structure du code source}

\subsection*{Répertoires principaux}

\begin{itemize}[label=\color{accentTeal}\faIcon{circle-dot}, leftmargin=*, itemsep=0.6em]
  \item \texttt{src/agents/} : Contient les agents métier (\texttt{extractor.py} avec \texttt{SecurityAgent})
  \item \texttt{src/utils/} : Modules utilitaires (\texttt{parsers.py}, \texttt{shared\_memory.py})
  \item \texttt{data/raw/} : Fichiers sources temporaires (logs, trafic, alertes)
  \item \texttt{data/processed/} : Canaux de mémoire partagée (JSON)
\end{itemize}

\subsection*{Module \texttt{parsers.py}}

\begin{tcolorbox}[
  enhanced,
  colback=white,
  colframe=accentTeal,
  left=6pt,
  right=6pt,
  top=6pt,
  bottom=6pt,
  boxrule=1pt,
  arc=4pt,
  before skip=8pt,
  after skip=8pt
]
\textbf{Fonctionnalités} :
\begin{itemize}[label=\color{primaryBlue}\faIcon{chevron-right}, leftmargin=*]
  \item \texttt{parse\_syslog\_line()} : Analyse syntaxique RFC~3164 avec extraction host/process/message
  \item \texttt{parse\_tshark\_json()} : Parsing du format EK/Elasticsearch pour trafic réseau
  \item \texttt{generate\_deterministic\_id()} : Génération SHA-256 d'ID sans LLM
  \item \texttt{normalize\_severity()} : Conversion vers score 0--1
\end{itemize}
\end{tcolorbox}

\subsection*{Module \texttt{extractor.py}}

L'agent principal \texttt{SecurityAgent} implémente :

\begin{enumerate}[label=\color{accentTeal}\faIcon{hashtag}, leftmargin=*, itemsep=0.5em]
  \item \textbf{Ingestion asynchrone} : Threads dédiés pour syslog (UDP/TCP) et capture réseau (tshark)
  \item \textbf{Gestion d'erreurs explicite} : Aucun échec silencieux -- tout problème est loggué avec \texttt{[❌ ERROR]} ou \texttt{[⚠️ WARNING]}
  \item \textbf{Arrêt gracieux} : Gestion \texttt{SIGINT}/\texttt{SIGTERM} avec nettoyage des sockets et des queues
  \item \textbf{Détection de pertinence} : Mots-clés de sécurité et scoring pour filtrage automatique
\end{enumerate}

\section*{Sources supportées}

\begin{itemize}[label=\color{accentTeal}\faIcon{server}, leftmargin=*, itemsep=0.8em]
  \item \textbf{Syslog (UDP/TCP)} : Port par défaut 1514, format RFC~3164
  \item \textbf{Capture réseau} : \texttt{tshark} avec format EK (\texttt{-T ek}) ou JSON classique
  \item \textbf{Alertes SIEM} : Fichier JSON/NDJSON depuis \texttt{data/raw/siem\_alerts.json}
\end{itemize}

\section*{Utilisation}

\subsection*{Mode daemon (temps réel)}

\begin{lstlisting}
python -m src.agents.extractor --daemon --syslog-port 1514
\end{lstlisting}

\subsection*{Mode batch (analyse de fichiers)}

\begin{lstlisting}
python -m src.agents.extractor --cleanup
\end{lstlisting}

\end{document}